% GENERATED FILE. Edit canonical lessons, then run npm run generate:books.
% canonical-source-hash: fnv1a64:3f55a808dcc1c732
% canonical-lessons: UR-C496-varkshap, UR-C496-shikanja, UR-C496-surakh, UR-C496-chhedna, UR-C496-regmal

\chapter{The Workshop}
\label{ch:ur-the-workshop}
\hlchaptermodality{496}

\begin{chapteropening}
\textbf{By the end of this chapter:} \emph{I can name things in a workshop and say what they do.}

Complete the last lesson of chapter 496: I can name things in a workshop and say what they do.
\end{chapteropening}

\section[\texorpdfstring{\ur{ورکشاپ} (varkshāp)}{ورکشاپ (varkshāp)}]{\ur{ورکشاپ} (varkshāp) — a workshop}

\label{lesson:UR-C496-varkshap}



\begin{quote}
Before the new one: say the Urdu for a contractor, then the Urdu for to use up (time or effort).
\end{quote}

\subsection*{You'll want to know: \ur{ورکشاپ}}
\textbf{\ur{ورکشاپ}} — \emph{varkshāp} — \textquotedblleft{}a workshop\textquotedblright{}.

\begin{quote}
\ur{ہماری} \ur{ورکشاپ} \ur{اسٹیشن} \ur{کے} \ur{قریب} \ur{ہے۔}
\end{quote}

\begin{quote}
\emph{hamārī varkshāp isṭeshan ke qarīb hai.}
\end{quote}

\begin{quote}
Our workshop is near the station.
\end{quote}

Feminine.

\begin{cousinweb}
\textbf{\ur{ورکشاپ}} (\emph{varkshāp}) is borrowed from English \emph{workshop}.
\end{cousinweb}

\begin{grammarlens}[title={What you've built}]
One more word for describing what you lived and felt.
\end{grammarlens}

\subsection*{Your turn}
\begin{itemize}
\raggedright
  \item \emph{Say it:} \emph{varkshāp}
  \item \emph{Say it:} \emph{varkshāp}, once more
  \item \emph{Say it:} \emph{hamārī varkshāp isṭeshan ke qarīb hai.} — Our workshop is near the station.
  \item \emph{Recall:} say the Urdu for a contractor, then the Urdu for to use up (time or effort), then say \emph{varkshāp} again
\end{itemize}

\begin{tcolorbox}[breakable,colback=teal!4,colframe=teal!35!black,title={Before you move on}]
What does \ur{ورکشاپ} mean? (\textquotedblleft{}A workshop\textquotedblright{}.) Say it once more.
\end{tcolorbox}
\section[\texorpdfstring{\ur{شکنجہ} (shikanjā)}{شکنجہ (shikanjā)}]{\ur{شکنجہ} (shikanjā) — a clamp; a vice}

\label{lesson:UR-C496-shikanja}



\begin{quote}
Before the new one: say the Urdu for to use up (time or effort), then the Urdu for a workshop.
\end{quote}

\subsection*{You'll want to know: \ur{شکنجہ}}
\textbf{\ur{شکنجہ}} — \emph{shikanjā} — \textquotedblleft{}a clamp; a vice\textquotedblright{}.

\begin{quote}
\ur{شکنجہ} \ur{میز} \ur{کے} \ur{کونے} \ur{پر} \ur{ہے۔}
\end{quote}

\begin{quote}
\emph{shikanjā mez ke kone par hai.}
\end{quote}

\begin{quote}
The clamp is on the corner of the table.
\end{quote}

Masculine.

\begin{cousinweb}
\textbf{\ur{شکنجہ}} (\emph{shikanjā}) is usually traced to Persian \emph{shikanja}, \textquotedblleft{}a press, a clamp\textquotedblright{}.
\end{cousinweb}

\begin{grammarlens}[title={What you've built}]
One more word for describing what you lived and felt.
\end{grammarlens}

\subsection*{Your turn}
\begin{itemize}
\raggedright
  \item \emph{Say it:} \emph{shikanjā}
  \item \emph{Say it:} \emph{shikanjā}, once more
  \item \emph{Say it:} \emph{shikanjā mez ke kone par hai.} — The clamp is on the corner of the table.
  \item \emph{Recall:} say the Urdu for to use up (time or effort), then the Urdu for a workshop, then say \emph{shikanjā} again
\end{itemize}

\begin{tcolorbox}[breakable,colback=teal!4,colframe=teal!35!black,title={Before you move on}]
What does \ur{شکنجہ} mean? (\textquotedblleft{}A clamp; a vice\textquotedblright{}.) Say it once more.
\end{tcolorbox}
\section[\texorpdfstring{\ur{سوراخ} (surākh)}{سوراخ (surākh)}]{\ur{سوراخ} (surākh) — a hole}

\label{lesson:UR-C496-surakh}



\begin{quote}
Before the new one: say the Urdu for a workshop, then the Urdu for a clamp.
\end{quote}

\subsection*{You'll want to know: \ur{سوراخ}}
\textbf{\ur{سوراخ}} — \emph{surākh} — \textquotedblleft{}a hole\textquotedblright{}.

\begin{quote}
\ur{دیوار} \ur{میں} \ur{ایک} \ur{چھوٹا} \ur{سوراخ} \ur{ہے۔}
\end{quote}

\begin{quote}
\emph{dīvār meṅ ek chhoṭā surākh hai.}
\end{quote}

\begin{quote}
There is a small hole in the wall.
\end{quote}

Masculine.

\begin{cousinweb}
\textbf{\ur{سوراخ}} (\emph{surākh}) is usually traced to Persian \emph{sūrākh}, \textquotedblleft{}a hole, an opening\textquotedblright{}.
\end{cousinweb}

\begin{grammarlens}[title={What you've built}]
One more word for describing what you lived and felt.
\end{grammarlens}

\subsection*{Your turn}
\begin{itemize}
\raggedright
  \item \emph{Say it:} \emph{surākh}
  \item \emph{Say it:} \emph{surākh}, once more
  \item \emph{Say it:} \emph{dīvār meṅ ek chhoṭā surākh hai.} — There is a small hole in the wall.
  \item \emph{Recall:} say the Urdu for a workshop, then the Urdu for a clamp, then say \emph{surākh} again
\end{itemize}

\begin{tcolorbox}[breakable,colback=teal!4,colframe=teal!35!black,title={Before you move on}]
What does \ur{سوراخ} mean? (\textquotedblleft{}A hole\textquotedblright{}.) Say it once more.
\end{tcolorbox}
\section[\texorpdfstring{\ur{چھیدنا} (chhednā)}{چھیدنا (chhednā)}]{\ur{چھیدنا} (chhednā) — to bore (a hole); to pierce}

\label{lesson:UR-C496-chhedna}



\begin{quote}
Before the new one: say the Urdu for a clamp, then the Urdu for a hole.
\end{quote}

\subsection*{You'll want to know: \ur{چھیدنا}}
\textbf{\ur{چھیدنا}} — \emph{chhednā} — \textquotedblleft{}to bore (a hole); to pierce\textquotedblright{}.

\begin{quote}
\ur{مستری} \ur{نے} \ur{دیوار} \ur{میں} \ur{دو} \ur{سوراخ} \ur{چھیدے۔}
\end{quote}

\begin{quote}
\emph{mistrī ne dīvār meṅ do surākh chhede.}
\end{quote}

\begin{quote}
The mechanic bored two holes in the wall.
\end{quote}

Past, with \ur{نے} (he or she …): \ur{اس} \ur{نے} \ur{دروازے} \ur{میں} \ur{سوراخ} \ur{چھیدا}, \emph{us ne darvāze meṅ surākh chhedā}.

\begin{cousinweb}
\textbf{\ur{چھیدنا}} (\emph{chhednā}) is inherited: CDIAL lists the Hindi cousin \emph{chednā}, \textquotedblleft{}to bore\textquotedblright{}, under Sanskrit \emph{chidra}, \textquotedblleft{}pierced; a hole\textquotedblright{} (CDIAL 5043).
\end{cousinweb}

\begin{grammarlens}[title={What you've built}]
One more word for describing what you lived and felt.
\end{grammarlens}

\subsection*{Your turn}
\begin{itemize}
\raggedright
  \item \emph{Say it:} \emph{chhednā}
  \item \emph{Say it:} \emph{chhednā}, once more
  \item \emph{Say it:} \emph{mistrī ne dīvār meṅ do surākh chhede.} — The mechanic bored two holes in the wall.
  \item \emph{Recall:} say the Urdu for a clamp, then the Urdu for a hole, then say \emph{chhednā} again
\end{itemize}

\begin{tcolorbox}[breakable,colback=teal!4,colframe=teal!35!black,title={Before you move on}]
What does \ur{چھیدنا} mean? (\textquotedblleft{}To bore (a hole); to pierce\textquotedblright{}.) Say it once more.
\end{tcolorbox}
\section[\texorpdfstring{\ur{ریگمال} (regmāl)}{ریگمال (regmāl)}]{\ur{ریگمال} (regmāl) — sandpaper}

\label{lesson:UR-C496-regmal}



\begin{quote}
Before the new one: say the Urdu for a hole, then the Urdu for to bore (a hole).
\end{quote}

\subsection*{You'll want to know: \ur{ریگمال}}
\textbf{\ur{ریگمال}} — \emph{regmāl} — \textquotedblleft{}sandpaper\textquotedblright{}.

\begin{quote}
\ur{بڑھئی} \ur{نے} \ur{ریگمال} \ur{سے} \ur{میز} \ur{چکنی} \ur{کی۔}
\end{quote}

\begin{quote}
\emph{baṛhaī ne regmāl se mez chiknī kī.}
\end{quote}

\begin{quote}
The carpenter smoothed the table with sandpaper.
\end{quote}

Masculine.

\begin{cousinweb}
\textbf{\ur{ریگمال}} (\emph{regmāl}) is usually traced to Persian \emph{rēgmāl}, built from \emph{rēg}, \textquotedblleft{}sand\textquotedblright{}, and \emph{māl-}, \textquotedblleft{}to rub\textquotedblright{}.
\end{cousinweb}

\begin{grammarlens}[title={What you've built}]
One more word for describing what you lived and felt.
\end{grammarlens}

\subsection*{Your turn}
\begin{itemize}
\raggedright
  \item \emph{Say it:} \emph{regmāl}
  \item \emph{Say it:} \emph{regmāl}, once more
  \item \emph{Say it:} \emph{baṛhaī ne regmāl se mez chiknī kī.} — The carpenter smoothed the table with sandpaper.
  \item \emph{Recall:} say the Urdu for a hole, then the Urdu for to bore (a hole), then say \emph{regmāl} again
\end{itemize}

\begin{tcolorbox}[breakable,colback=teal!4,colframe=teal!35!black,title={Before you move on}]
What does \ur{ریگمال} mean? (\textquotedblleft{}Sandpaper\textquotedblright{}.) Say it once more.
\end{tcolorbox}
